\chapter{魔術の残骸と、冷徹なる計算機}
\epigraph{``The brain is not just a computer, it's a learning machine''.\\（脳は単なるコンピュータではない。自ら学習する機械なのだ。）}{Dr.~Richard Bandler}
NLP（神経言語プログラミング）が「再現性のない疑似科学」と断じられて久しい。アカデミアの書棚から追放され、自己啓発セミナーの商材として消費される現状を見て、「やはりインチキだったのか」と失望した者も多いだろう。あるいは、ラポールやイエスセットといった表層的なテクニックの一片を切り取り、「効果がある」と安易に信じ込んでいる者もいる。

だが、現場の真実を知る者は知っている。
私たちは、リチャード・バンドラー博士が放つ圧倒的な「魔術」をこの目で目撃してきた。重度の恐怖症が一瞬で霧散し、数十年来の頑なな信念がわずか数分の対話で解体され、再構築される。彼の指先ひとつ、声のトーンひとつで、人間の知覚空間そのものが歪められる。彼は間違いなく本物である。

しかし――なぜ私たちは、彼と同じスクリプトを読み上げ、同じ催眠言語パターンをなぞっても、一向に相手をトランスに落とせないのか。なぜ他人はこちらの指示に従わず、強固な現実の壁に跳ね返されるのか。

答えは残酷なほどシンプルだ。
私たちは彼が「見せている手品（表層の言語）」を表面的に真似ていただけで、彼がハックしていた「生体情報処理ハードウェアの物理法則」を何ひとつ理解していなかったからである。

NLPを神秘のベールから引き剥がし、真の工学として掌中に収めるためには、私たちが「心」や「無意識」に対して抱いてきた自明の思い込みを疑う、冷徹な批判的知性が必要となる。

NLPを神秘のベールから引き剥がし、真の工学として掌中に収めるためには、私たちが「心」や「無意識」に対して抱いてきた自明の思い込みはもちろんのこと、NLPという体系そのものが掲げる前提すら疑う冷徹な批判的知性が必要となる。

種明かしを始める前に、まずはこの生体コンピュータの動作原理を記述するための、最低限の数学的・情報論的プロトコルを定義しておこう。

\section{ベイズ統計学の基礎数理}
\epigraph{``For 150 years, Bayes was taboo... dismissed as subjective and unscientific.''\\(ベイズ統計は、150年ものあいだタブーとされ、異端審問にかけられた。主観を計算式に入れるなど、科学の裏切りだと見なされたからだ。)}{Sharon Bertsch McGrayne}
ベイズ統計学とは、新しいデータや証拠（エビデンス）を獲得するたびに、自らの見立て（確信度）を柔軟に更新していく実践的な確率論である。
無限回の繰り返し試行を前提とする従来の頻度主義統計学とは異なり、ベイズ統計学は「一度きりの事象」に対しても主観確率を割り振って評価できる。これにより、情報が極めて乏しい初期段階であっても事前知識をもとに暫定的な推論を立ち上げることができる。

\subsection{事象の確率と条件付き確率}
ベイズ推論を理解するための基礎となるのが「条件付き確率（Conditional Probability）」である。

確率が正である事象 $B$（すなわち $P(B) > 0$）がすでに観測された条件下で、事象 $A$ が生起する条件付き確率は次式で定義される。
\begin{equation}
P(A \mid B) = \frac{P(A \cap B)}{P(B)}
\label{eq:cond_prob}
\end{equation}

ここで $\Probability{\eventA \cap \eventB}$ は事象 $\eventA$ と事象 $\eventB$ が同時に生起する同時確率を表す。この定義式を変形すると、確率の乗法定理が得られる。
\begin{equation}
P(A \cap B) = P(A \mid B) P(B) = P(B \mid A) P(A)
\label{eq:mult_rule}
\end{equation}
\eqref{eq:mult_rule} 式は、事象 $A$ と $B$ が同時に生起するプロセスを「$B$ を前提とした $A$ の生起」と捉えても、「$A$ を前提とした $B$ の生起」と捉えても数学的に同値であることを示している。

\subsection{ベイズの定理の導出}
乗法定理 \eqref{eq:mult_rule} の両辺を $P(B)$（ただし $P(B) > 0$）で除すことにより、基礎的なベイズの定理（Bayes' Theorem）が導かれる。
\begin{equation}
P(A \mid B) = \frac{P(B \mid A) P(A)}{P(B)}
\label{eq:bayes_basic}
\end{equation}
この関係式は、ある方向の条件付き確率 $P(B \mid A)$ が既知であるとき、その逆方向の条件付き確率 $P(A \mid B)$ を計算できることを示している。

\subsection{全確率の定理を用いた一般化}
起こりうるすべての状態が、互いに排他で網羅的な $K$ 個の事象の分割 $\{H_1, H_2, \dots, H_K\}$ によって構成されているとする。
すなわち、
\begin{equation}
    \bigcup_{j=1}^K H_j = \Omega_{\mathrm{event}}, \quad H_i \cap H_j = \emptyset \quad (i \neq j)
\end{equation}
このとき、任意の事象 $D$（観測データ）の確率 $P(D)$ は「全確率の定理」により次のように展開できる。
\begin{equation}
P(D) = \sum_{j=1}^K P(D \cap H_j) = \sum_{j=1}^K P(D \mid H_j) P(H_j)
\label{eq:law_of_total_prob}
\end{equation}
これを \eqref{eq:bayes_basic} 式の分母に代入することで、一般化されたベイズ方程式が得られる。
\begin{equation}
P(H_k \mid D) = \frac{P(D \mid H_k) P(H_k)}{\displaystyle\sum_{j=1}^K P(D \mid H_j) P(H_j)}
\label{eq:gen_bayes}
\end{equation}

この方程式を構成する各項は、認知および統計的推論において次のように定義される。
\begin{itemize}
  \item \textbf{事前確率（Prior Probability）} $P(H_k)$：データ $D$ を得る前に、仮説 $H_k$ が正しいと見積もられている確信度。
  \item \textbf{尤度（Likelihood）} $P(D \mid H_k)$：仮説 $H_k$ が真である場合に、データ $D$ が観測される確率。
  \item \textbf{周辺尤度／正規化定数（Marginal Likelihood / Evidence）} $P(D) = \sum_{j=1}^K P(D \mid H_j) P(H_j)$：すべての仮説のもとでデータ $D$ が観測される総合確率。
  \item \textbf{事後確率（Posterior Probability）} $P(H_k \mid D)$：データ $D$ を観測した後に更新された、仮説 $H_k$ に対する確信度。
\end{itemize}

分母 $P(D)$ は各仮説 $H_k$ に依存しない共通の正規化定数であるため、ベイズ方程式は本質的に次の比例関係へと帰着する。
\begin{equation}
P(H_k \mid D) \propto P(D \mid H_k) \times P(H_k)
\label{eq:proportional}
\end{equation}
すなわち、「\textbf{事後確率 $\propto$ 尤度 $\times$ 事前確率}」である。

\subsection{ベイズ更新（逐次更新）のメカニズム}
ベイズ推論の真価は、時系列で入力されるデータに対し、前回の事後確率を次の事前確率へと繰り越しながらリアルタイムに確信度を再計算できる「逐次性」にある。

1回目のデータ $D_1$ を観測したとき、第1段階の事後確率は次のように得られる。

\begin{equation}
P(H_k \mid D_1) = \frac{P(D_1 \mid H_k) P(H_k)}{P(D_1)}
\label{eq:step1}
\end{equation}

続いて独立な2つ目のデータ $D_2$ が追加されたとき、直前の事後確率 $P(H_k \mid D_1)$ を\textbf{新たな事前確率}として採用する。
仮説 $H_k$ のもとでデータ群が条件付き独立（すなわち $P(D_2 \mid H_k, D_1) = P(D_2 \mid H_k)$）である場合、2段階目の更新式は次のように書ける。
\begin{equation}
P(H_k \mid D_1, D_2) = \frac{P(D_2 \mid H_k) P(H_k \mid D_1)}{\displaystyle\sum_{j=1}^K P(D_2 \mid H_j) P(H_j \mid D_1)}
\label{eq:step2_final}
\end{equation}

この逐次更新の結果は、すべての観測データ $\{D_1, D_2\}$ を一度にまとめて一括推論した結果と厳密に一致する。
\begin{align}
P(H_k \mid D_1, D_2) 
&\propto P(D_2 \mid H_k) \cdot P(H_k \mid D_1) \notag \\
&= P(D_2 \mid H_k) \cdot \left( \frac{P(D_1 \mid H_k) P(H_k)}{P(D_1)} \right) \notag \\
&\propto P(D_1 \mid H_k) P(D_2 \mid H_k) P(H_k) \notag \\
&= P(D_1, D_2 \mid H_k) P(H_k)
\end{align}

\section{モンティ・ホール問題：人間の直感の破綻}
\epigraph{``Our brains are just not wired for probability problems.''\\
（人間の脳は、確率計算ができるようには最初から配線されていないのだ。）}{Paul Erd\H{o}s}

人間の神経系がいかに厳密なベイズ計算から逸脱した「歪んだ推論機」であるかを浮き彫りにした歴史的事件が、モンティ・ホール問題である。

1963年にアメリカで放送が開始された大人気番組『Let's Make a Deal』において、司会者モンティ・ホールは次のようなゲームを提供していた。
\begin{quote}
3つのドア（1番、2番、3番）があり、1つの後ろには「新車（当たり）」、残り2つの後ろには「ヤギ（ハズレ）」がいる。
挑戦者が1つのドアを選んだ後、正解を知っている司会者モンティが、残りのドアから「必ずヤギがいるハズレのドア」を1つ開けて見せる。
その上で、モンティは挑戦者にこう提案する。
「ドアを変更しますか？ それとも最初の選択のままにしますか？」
\end{quote}

\subsection{全米を巻き込んだ大論争}
1990年、ギネス世界最高IQ保持者と称されていたマリリン・ボス・サヴァントが雑誌コラムで「\textbf{絶対にドアを変えるべき。勝率は2倍（$2/3$）になる}」と回答した。
これに対し、全米から約1万通の猛抗議が殺到した。その中には約1,000人の数学者・博士号保持者が含まれており、「ドアが2枚になったのだから確率は半々（$1/2$）に決まっている」「プロの数学者として恥を知れ」と激しくバッシングした。

しかし、コンピュータシミュレーションや厳密な数理検証の結果、勝率はドアを変えることで正確に $2/3$（約66.7\%）になることが実証された。批判した数学者たちは公式に謝罪へ追い込まれた。

\subsection{ベイズの定理による厳密解}
挑戦者が最初に「ドア1」を選び、モンティがハズレの「ドア3」を開けた状況をベイズの定理で計算する。

車がある位置の仮説を $C_i$（$i \in \{1, 2, 3\}$）、モンティがドア3を開ける事象を $M_3$ とする。
初期状態の事前確率は等配分である。
\begin{equation}
P(C_1) = P(C_2) = P(C_3) = \frac{1}{3}
\end{equation}

次に、各仮説下でモンティがドア3を開ける尤度 $P(M_3 \mid C_i)$ を評価する。モンティは「挑戦者の選んだドア」と「車のあるドア」を決して開けない。
\begin{enumerate}[label=(\roman*)]
  \item $C_1$（車がドア1にある）：残るドア2とドア3は共にハズレであるため、ランダムに開ける。
  \begin{equation}
  P(M_3 \mid C_1) = \frac{1}{2}
  \end{equation}
  \item $C_2$（車がドア2にある）：車があるドア2は開けられないため、必然的にドア3を開けるしかない。
  \begin{equation}
  P(M_3 \mid C_2) = 1
  \end{equation}
  \item $C_3$（車がドア3にある）：車があるドア3を開けることはルール上あり得ない。
  \begin{equation}
  P(M_3 \mid C_3) = 0
  \end{equation}
\end{enumerate}

全確率の定理より、モンティがドア3を開ける周辺尤度は次のようになる。
\begin{align}
P(M_3) &= \sum_{i=1}^3 P(C_i) P(M_3 \mid C_i) \notag \\
       &= \left(\frac{1}{3} \times \frac{1}{2}\right) + \left(\frac{1}{3} \times 1\right) + \left(\frac{1}{3} \times 0\right) = \frac{1}{2}
\end{align}

したがって、ベイズの定理による事後確率は以下の通りとなる。
\begin{align}
P(C_1 \mid M_3) &= \frac{P(C_1) P(M_3 \mid C_1)}{P(M_3)} = \frac{\frac{1}{3} \times \frac{1}{2}}{\frac{1}{2}} = \frac{1}{3} \\
P(C_2 \mid M_3) &= \frac{P(C_2) P(M_3 \mid C_2)}{P(M_3)} = \frac{\frac{1}{3} \times 1}{\frac{1}{2}} = \frac{2}{3}
\end{align}

ドアを変えることで勝率は正確に2倍に跳ね上がる。
多くの学者すら錯誤した原因は、「情報開示の非対称性（車がドア2にある場合、モンティの選択肢がドア3に狭まるという拘束条件）」を無視し、目の前の2枚のドアという対称な静的スナップショットに認知が囚われた点にある。

人間の直感は、数学的に完全なベイズ推論機ではない。生存のために都合よくバイアスを組み込んだ\textbf{「不完全な疑似ベイズ推論機」}にすぎないのだ。そして、学習とは数学的ベイズ更新に似ているが異なる\textbf{「疑似ベイズ更新」}なのだ。


